% Track e7 -- Classical readings: astronomy and cosmology.
% Verification: investigations/emergence/verify/e7_classical.py (claims e7-1 ... e7-27).
% New references: sections/bib_e7.tex.
\section{Classical readings: astronomy and cosmology}
\label{e7:sec:main}

This section reads the results of \cite{RE} and of the preceding sections in the sense of classical physics, astronomy and cosmology. Five
classes are kept apart. A statement is \emph{inherited} when it is ordinary physics, true in any notation, and cited from the literature. It is
a \emph{reading} (interpretive, in the classes of this note) when it identifies a quantity of the arena with an observed one; we say what the
identification needs. It is a \emph{constraint} when, if a reading is adopted, an observation bounds a parameter. The constraint is then fixed
arithmetic on cited inputs, conditional on the reading. It is \emph{obstructed} when observation or inherited physics excludes the reading, and
\emph{open} otherwise. Mathematical statements carry the classes fixed and restated as before.

One fact governs the whole section. The arena is scale-free: the rate $\kappa=c/\ell$ of the interior flow has no fixed length $\ell$, the
identification $mc^2=\hbar\kappa$ is added, and the multiscale ratio of Corollary~17.15 of \cite{ES17} is an input (Theorem~4.12,
Propositions~4.8 and~4.9 of \cite{RE}). No number of astronomy or cosmology follows from the arena alone. Every number below is either
inherited arithmetic on observed inputs or a bound on a parameter of a reading.

The answers, with their classes, are these. (1) The rest energy $u_0$ of a field point is the energy density in the frames where $\bE\parallel\Btext$
(inherited). Its stress is a tension $u_0$ along the common axis and a pressure $u_0$ across it, and every isotropic ensemble of field points
gravitates as radiation, $w=\tfrac13$, whatever its $u_0$. So the reading ``field rest energy is dark matter'' is obstructed
(\cref{e7:sec:point}). (2) For isotropic thermal radiation at a point the ratio $E_0/E_6$ has density $2t$ on $[0,1]$, independent of the
energy, so $\langle E_0\rangle/\langle E_6\rangle=\tfrac23$ (fixed). Every plane wave is null, but the random superposition is massive almost
everywhere; ``the CMB is locally massive'' is a reading (\cref{e7:sec:cmb}). (3) The pair invariant of two null fields of \cite{RE} vanishes
on the same configurations as the Breit--Wheeler invariant. The thresholds of the gamma-ray horizon need $m_e$, which the arena does not
supply. The vacuum pair rate depends on the canonical electric field $\epsilon$, not on $E_0=\half(\epsilon^2+\beta^2)$, so magnetar fields above
$B_Q$ create no pairs from the vacuum (inherited; \cref{e7:sec:pairs}). (4) Two six-momenta at mutual rest with temporal angle $\vartheta$ have
$M^2=m_1^2+m_2^2+2m_1m_2\cos\vartheta$. In three times ``uphill'' decays are kinematically open, and one clock closes them (inherited, with an
exact example). Mass additivity bounds the temporal angle of the neutron and $^{28}$Si to $9\times10^{-5}$~rad, if the reading is adopted
(\cref{e7:sec:sixmomenta}). (5) A duality rotation of a null field is a rotation about its propagation axis (fixed). If the interior flow acted
on light, CMB polarisation would bound its rate by $2.5\times10^{-20}\,{\rm s^{-1}}$, $10^{40.5}$ below $m_ec^2/\hbar$; that reading is
obstructed (\cref{e7:sec:rate}). (6) The three-manifold time-odd invariant of \cite{RE} parallels the Jarlskog invariant only as a reading. A
temporal $b$-field acting on the clocks of matter is bounded by the constancy of $m_p/m_e$ and by the equivalence principle, as conditional
constraints (\cref{e7:sec:timeodd}). \Cref{e7:sec:summary} collects the claims (\cref{e7:tab:claims}). Identifiers e7-$n$ refer to the
\texttt{PASS} lines of \texttt{verify/e7\_classical.py}; physical constants are CODATA values from \texttt{scipy.constants}.

\paragraph{Notation.} SI units. The frame variable is $\bB=-\Btext$, with $\Btext$ the textbook field (Section~1.4 of \cite{ES17}), $z_c=c\bB+i\bE$,
$q_c=z_c\cdot z_c$, $u=\tfrac{\varepsilon_0}2|z_c|^2$, $u_0=\tfrac{\varepsilon_0}2|q_c|$, $\mathbf g=\bE\times\Btext/(\mu_0c^2)$ and $\rho_0=u_0/c^2$. Then
$u^2=c^2|\mathbf g|^2+u_0^2$ (Proposition~2.8 of \cite{RE}; restated, e7-2). In Heaviside--Lorentz units with $c=1$, $z=\bB+i\bE$, $E_6=\half|z|^2$,
$E_0=\half|q|$, $q=z\cdot z=|\Btext|^2-|\bE|^2-2i\,\bE\cdot\Btext$. The realisation is that of \cite{ES17}, and $G$ acts on $z$ as $i$ (e7-1).

\subsection{The rest energy of a field point in the sky}
\label{e7:sec:point}

The minimum of $u$ over frames is $u_0$. It is attained exactly in the frames where the Poynting vector vanishes, and there $\bE\parallel\Btext$
or one of them vanishes; for $q\neq0$ such frames exist (Theorem~2.6 of \cite{RE}). This is inherited: the classification of a field at a point
\cite{Synge56} and the problem of the frame with parallel fields in \S25 of \cite{LandauLifshitz2}. Let
$T^{00}=u$, $T^{0i}=(\bE\times\Btext)_i$, $T^{ij}=-(E_iE_j+B_{{\rm t}i}B_{{\rm t}j})+u\,\delta_{ij}$ be the stress--energy tensor of the field
(Heaviside--Lorentz, $c=1$).

\begin{proposition}[The stress of a field point]
\label{e7:prop:stress}
\begin{enumerate}[label=(\roman*),leftmargin=2.2em]
\item $T^\mu{}_\mu=0$ and $T^\mu{}_\lambda T^\lambda{}_\nu=E_0^2\,\delta^\mu{}_\nu$.
\item If $\bE=\epsilon\,\mathbf n$ and $\Btext=\beta\,\mathbf n$ with $|\mathbf n|=1$, then $T^{00}=u_0$, $T^{0i}=0$ and $T^{ij}=u_0(\delta_{ij}-2n_in_j)$:
a tension $u_0$ along $\mathbf n$ and a pressure $u_0$ in the two directions across it.
\item In every ensemble of field points whose mean stress is isotropic in some frame, $\langle T^{ij}\rangle=p\,\delta^{ij}$ with
$p=\tfrac13\langle T^{00}\rangle$, whatever the distribution of $u_0$. For field points at rest with isotropically distributed axes,
$p=\tfrac13\langle u_0\rangle$.
\item In a Friedmann--Lema\^itre universe a component with $p=w\rho c^2$ obeys $\dot\rho=-3(1+w)\rho\,\dot a/a$. So $\rho\propto a^{-4}$ for every
such ensemble of free fields, against $a^{-3}$ for dust.
\end{enumerate}
\end{proposition}

\begin{proof}
(i) $\sum_iT^{ii}=-(|\bE|^2+|\Btext|^2)+3u=u$, so $T^\mu{}_\mu=T^{00}-\sum_iT^{ii}=0$; the square is the Rainich identity (Proposition~2.5 of
\cite{RE}). (ii) Substitute; $u=\half(\epsilon^2+\beta^2)=u_0$ since $|q|=\epsilon^2+\beta^2$ when $\bE\parallel\Btext$, and $\bE\times\Btext=0$.
(iii) Tracelessness gives $\langle T^{00}\rangle=\sum_i\langle T^{ii}\rangle=3p$. In the second case $\langle n_in_j\rangle=\tfrac13\delta_{ij}$, so
$\langle T^{ij}\rangle=u_0(1-\tfrac23)\delta_{ij}$. (iv) Integrate the continuity equation.
\end{proof}

Inherited (\S\S33 and~35 of \cite{LandauLifshitz2}; \cite{Rainich1925} via \cite{RE}; \cite{WeinbergCosmo08}), checked exactly (e7-3, e7-4). A
bounded static configuration together with its sources has total mass $U/c^2$, with the field stresses balanced by those of the sources (von Laue's
theorem \cite{Laue1911}; Proposition~2.9 of \cite{RE}). Field energy then gravitates as part of the mass of a matter system, and the matter fixes
its share.

\begin{corollary}[Field rest energy is not dark matter]
\label{e7:cor:dm}
Reading $\rho_0=u_0/c^2$ as a rest-mass density does not make field rest energy behave as matter in cosmology. The reading ``field rest energy is
dark matter'' is obstructed by inherited physics.
\end{corollary}

\begin{proof}
What gravitates is $T^{\mu\nu}$, and the reading leaves it unchanged. The dark matter of the standard model of cosmology is pressureless
since before recombination, with $\Omega_ch^2=0.120\pm0.001$ \cite{Planck18VI}. By \cref{e7:prop:stress}(iii)--(iv) every isotropic ensemble
of free field points has $w=\tfrac13$ and redshifts as $a^{-4}$, whatever $u_0$. A field point with $u_0>0$ is at rest only in the sense of
\cref{e7:prop:stress}(ii): its stress is not that of dust, and a free field propagates at $c$. The only configurations whose field energy
weighs like rest mass are bound ones, by von Laue's theorem, and there the mass belongs to the matter system. The amounts are also small: the
mean cold-dark-matter energy density is $1.26\,{\rm keV\,cm^{-3}}$, of which the local galactic magnetic energy density is $7\times10^{-4}$ and the
mean rest energy of the CMB (\cref{e7:sec:cmb}) $1.4\times10^{-4}$ (e7-26). The cause of the obstruction is the tracelessness, not the deficit.
\end{proof}

The fields of astronomy are, to a good approximation, magnetic in the frame of the conducting matter, where $\bE\approx0$, so $u=u_0=B^2/(2\mu_0)$
and $\rho_0=B^2/(2\mu_0c^2)$ there. \Cref{e7:tab:fields} gives the arithmetic (e7-5) for representative inputs: the geomagnetic field at the surface
(about $50\,\mu$T; total intensity maps of \cite{Alken21}), a sunspot umbra ($0.3$~T \cite{Solanki03}), a strongly magnetic white dwarf (isolated
magnetic white dwarfs reach about $10^5$~T \cite{Ferrario15}), the surface of a radio pulsar (about $10^8$~T \cite{LorimerKramer05}), a magnetar
($10^{10}$--$10^{11}$~T \cite{KaspiBeloborodov17}) and the galactic field (a few microgauss, about $6\,\mu$G in total near the Sun
\cite{Beck16}).

\begin{table}[ht]
\centering
\small
\begin{tabular}{@{}lccc@{}}
\toprule
field & $B$ (T) & $u=u_0$ (J\,m$^{-3}$) & $\rho_0$ (kg\,m$^{-3}$)\\
\midrule
geomagnetic, surface & $5\times10^{-5}$ & $9.9\times10^{-4}$ & $1.1\times10^{-20}$\\
sunspot umbra & $0.3$ & $3.6\times10^{4}$ & $4.0\times10^{-13}$\\
magnetic white dwarf & $10^{5}$ & $4.0\times10^{15}$ & $4.4\times10^{-2}$\\
radio pulsar, surface & $10^{8}$ & $4.0\times10^{21}$ & $4.4\times10^{4}$\\
magnetar & $10^{10}$--$10^{11}$ & $4.0\times10^{25}$--$4.0\times10^{27}$ & $4.4\times10^{8}$--$4.4\times10^{10}$\\
galactic, $3$--$6\,\mu$G & $3$--$6\times10^{-10}$ & $3.6$--$14\times10^{-14}$ & $4.0$--$16\times10^{-31}$\\
\midrule
CMB, $T_0=2.72548$~K & --- & $4.17\times10^{-14}$ ($u$) & $3.10\times10^{-31}$ ($\langle\rho_0\rangle$)\\
\bottomrule
\end{tabular}
\caption{Energy and rest-energy densities of astronomical fields in the frame where $\bE\approx0$ (e7-5), and of the CMB (e7-10). The last
column is $u_0/c^2$; reading it as a mass density is the reading of \cref{e7:cor:dm}.}
\label{e7:tab:fields}
\end{table}

The galactic field of $3$--$6\,\mu$G carries $0.22$--$0.89\,{\rm eV\,cm^{-3}}$. The CMB carries $aT_0^4=0.261\,{\rm eV\,cm^{-3}}$ ($T_0=2.72548$~K
\cite{Fixsen09}), and the cosmic rays of the local interstellar medium about $1\,{\rm eV\,cm^{-3}}$ \cite{Ferriere01,Draine11}. The three agree to a
factor of a few (e7-6). This rough equipartition of the magnetic, cosmic-ray and thermal energy densities is an inherited fact of the interstellar
medium \cite{Ferriere01}. The arena neither explains it nor uses it.

\subsection{The pointwise rest energy of isotropic radiation}
\label{e7:sec:cmb}

\begin{proposition}[Rest-energy fraction of a Gaussian field]
\label{e7:prop:gauss}
Let $z=\mathbf x+i\mathbf y$ with $\mathbf x,\mathbf y\in\R^3$ independent and $N(0,\sigma^2\id)$, and let $\lambda_1\ge\lambda_2$ be the
eigenvalues of the Gram matrix $\Gamma$ of $(\mathbf x,\mathbf y)$. Then $|q|=\lambda_1-\lambda_2$ and $|z|^2=\lambda_1+\lambda_2$. The ratio
$t=E_0/E_6=|z\cdot z|/|z|^2\in[0,1]$ has density $2t$ (distribution function $t^2$), is independent of $|z|^2$, and has mean $\tfrac23$ and
second moment $\half$. Hence $\langle E_0\rangle/\langle E_6\rangle=\tfrac23$, and $q\neq0$ with probability one.
\end{proposition}

\begin{proof}
$|q|^2=(\mathbf x\cdot\mathbf x-\mathbf y\cdot\mathbf y)^2+4(\mathbf x\cdot\mathbf y)^2=(\tr\Gamma)^2-4\det\Gamma=(\lambda_1-\lambda_2)^2$. $\Gamma$ is a real Wishart
matrix $W_2(3,\sigma^2\id)$, whose eigenvalues have the density $\propto(\lambda_1-\lambda_2)e^{-(\lambda_1+\lambda_2)/2\sigma^2}$ on
$\lambda_1>\lambda_2>0$; the general factor $\prod_i\lambda_i^{(n-p-1)/2}$ is $1$ at $n=3$, $p=2$ \cite{Muirhead82}. With $s=\lambda_1+\lambda_2$,
$t=(\lambda_1-\lambda_2)/s$ and Jacobian $s/2$ the density becomes $\propto s^2e^{-s/2\sigma^2}\cdot t$ on $s>0$, $0<t<1$. It factorises, the first
factor being the $\chi^2_6$ law of $|z|^2/\sigma^2$. So $t$ has density $2t$ and is independent of $s$, and
$\langle ts\rangle/\langle s\rangle=\langle t\rangle=\tfrac23$.
\end{proof}

Fixed: symbolically (e7-7); by Isserlis' theorem, $\langle|q|^2\rangle=24\sigma^4$, $\langle|z|^4\rangle=48\sigma^4$,
$\langle|q|^4\rangle=1920\sigma^8$ and $\langle|z|^8\rangle=5760\sigma^8$, whose ratios are the moments $\half$ and $\tfrac13$ of the density $2t$
(e7-8); and by Monte Carlo with $2\times10^6$ samples, with mean $0.6665$, Kolmogorov distance $7.5\times10^{-4}$ to $t^2$ and correlation
$-1.1\times10^{-3}$ between $t$ and $|z|^2$ (e7-9).

\emph{The model.} Thermal light is a circular complex Gaussian field, by the central limit theorem over many independent emitters
(Sections~4.2 and~4.3 of \cite{Goodman15}); this is inherited. For isotropic, unpolarised blackbody radiation at a point, isotropy gives
$\langle E_iE_j\rangle=\langle B_{{\rm t}i}B_{{\rm t}j}\rangle=\tfrac13u\,\delta_{ij}$ and $\langle E_iB_{{\rm t}j}\rangle=\tfrac13\langle\bE\cdot\Btext\rangle\delta_{ij}$,
and parity invariance gives $\langle\bE\cdot\Btext\rangle=0$. The six components of $z$ are therefore independent with a common variance, which
is the hypothesis of \cref{e7:prop:gauss}, and $\langle E_6\rangle=u$, $\langle E_0\rangle=\tfrac23u$. The distribution $2t$ is a prediction of the
Gaussian model. Any measurement of the joint statistics of $\bE$ and $\Btext$ at a point of thermal radiation can decide it; it is not a
prediction of the arena.

\emph{The reading.} Every plane wave is null, $q=0$ (Theorem~13.8 of \cite{ES17}). A random superposition has $q\neq0$ almost everywhere, with
mean rest-energy density $\tfrac23u$. To say that the CMB is ``locally massive'' is a reading. Its content is $\langle E_0\rangle=\tfrac23u$, and it
says nothing about gravitation, which sees the traceless $T^{\mu\nu}$ of \cref{e7:prop:stress}. The whole behaves differently. Radiation in a box at
rest in the CMB frame has $\mathbf P=0$ and $Mc^2=U$, while $\int E_0=\tfrac23U$ on average; this is the strict inequality $M\ge\int E_0$ of
Proposition~2.9(2) of \cite{RE}, strict because the local energy--momenta are not proportional. Numbers (e7-10): today
$u=aT_0^4=4.17\times10^{-14}\,{\rm J\,m^{-3}}$, $\langle u_0\rangle=2.78\times10^{-14}\,{\rm J\,m^{-3}}$ and
$\langle\rho_0\rangle=3.10\times10^{-31}\,{\rm kg\,m^{-3}}=3.6\times10^{-5}\rho_{\rm crit}$, with $\rho_{\rm crit}=3H_0^2/8\pi G$ for
$H_0=67.4\,{\rm km\,s^{-1}Mpc^{-1}}$ \cite{Planck18VI}. Like the whole CMB ($\Omega_\gamma=5.4\times10^{-5}$), this density redshifts as $a^{-4}$.

\subsection{Two from one in the sky: pair creation}
\label{e7:sec:pairs}

\emph{Two photons.} For null four-momenta $k_1,k_2$ with energies $E_1,E_2$ at angle $\vartheta$, $s=(k_1+k_2)^2=2E_1E_2(1-\cos\vartheta)$, and
$\gamma\gamma\to e^+e^-$ needs $s\ge4m_e^2c^4$ \cite{BreitWheeler34}. Head-on, $E_1E_2\ge(m_ec^2)^2$ (e7-11). On CMB photons of mean energy
$\pi^4kT_0/30\zeta(3)=6.3\times10^{-4}$~eV the threshold is $4.1\times10^{14}$~eV, and on photons of energy $kT_0$ it is $1.1\times10^{15}$~eV. On
optical and near-infrared extragalactic background light, $0.1$--$3$~eV, it is $0.09$--$2.6$~TeV. These are the edges of the gamma-ray horizon
\cite{Nikishov62,GouldSchreder67,DwekKrennrich13}; Gould and Schr\'eder predicted the cut-off of the cosmic photon spectrum above $10^{14}$~eV by
the $3$~K photons. The process has been observed with quasi-real, linearly polarised photons of ultraperipheral Au+Au collisions, with angular
distributions in accord with QED \cite{STAR21}. All of this is inherited.

\emph{Relation to the arena.} For two null fields $M_{\rm pair}=|z_1\cdot z_2|=(E_6^{(1)}E_6^{(2)})^{1/2}(1-\cos\vartheta_{12})$, with $\vartheta_{12}$ the
angle between the Poynting vectors (Theorem~8.6(v) of \cite{RE}; restated, e7-12). One null field does not become a massive pair under any
covariant duplication, and no closure dynamics fills an empty constituent (Corollary~8.3 of \cite{RE}; \cref{e2:cor:field}). The precise
relation is this. $M_{\rm pair}$ and $s$ vanish on the same configurations, the collinear ones, and are positive on all others: the selection
rule ``one photon makes no pair, two non-collinear photons can'' is shared. The functions differ. $M_{\rm pair}$ is linear in $1-\cos\vartheta$,
because field densities are quadratic in the fields; $s$ is linear in $1-\cos\vartheta$, so the photon pair mass $s^{1/2}$ goes as its square root.
The threshold $2m_ec^2$ needs $m_e$, which the arena does not supply (Theorem~4.12 of \cite{RE}). The gamma-ray horizon tests QED and the
background light, not the arena.

\emph{Constant fields.} If $q\neq0$ there are frames with $\bE=\epsilon\,\mathbf n$, $\Btext=\beta\,\mathbf n$ (\S25 of \cite{LandauLifshitz2}).

\begin{proposition}[Canonical fields]
\label{e7:prop:canon}
In such a frame, with $\epsilon\ge0$ and $\beta$ real,
\[
\epsilon^2=\half(|q|-\Re q),\qquad \beta^2=\half(|q|+\Re q),\qquad E_0=\half|q|=\half(\epsilon^2+\beta^2).
\]
The characteristic polynomial of $F^\mu{}_\nu$ is $\lambda^4+(|\Btext|^2-|\bE|^2)\lambda^2-(\bE\cdot\Btext)^2=(\lambda^2-\epsilon^2)(\lambda^2+\beta^2)$, with
eigenvalues $\pm\epsilon$ and $\pm i\beta$: $\epsilon=\epsilon_E$ and $|\beta|=\epsilon_M$ of Section~14.1 of \cite{ES17}. A purely magnetic $F$
generates rotations and a purely electric one boosts.
\end{proposition}

\begin{proof}
$q$ is invariant under the Lorentz group, which acts on $z$ through $SO(3,\C)$. In the canonical frame $\bE\cdot\Btext=\epsilon\beta$, so
$q=\beta^2-\epsilon^2-2i\epsilon\beta=(\beta-i\epsilon)^2$, $|q|=\epsilon^2+\beta^2$ and $\Re q=\beta^2-\epsilon^2$. The coefficients of the
characteristic polynomial are invariants and take these values. A real antisymmetric $F$ exponentiates to rotations, a symmetric one to boosts.
\end{proof}

Inherited, checked exactly and on $200$ random fields to $10^{-14}$ (e7-13). The vacuum decay rate in a constant field is
\[
w=\frac{e^2\epsilon\beta}{4\pi^2\hbar^2c}\sum_{n\ge1}\frac1n\coth\frac{n\pi\beta}{\epsilon}\,\exp\Bigl(-\frac{n\pi E_S}{\epsilon}\Bigr),\qquad
E_S=\frac{m_e^2c^3}{e\hbar}=1.32\times10^{18}\,{\rm V\,m^{-1}},
\]
in Heaviside--Lorentz units \cite{Schwinger51,HeisenbergEuler36}. It is exponential in $\epsilon$ and vanishes for $\epsilon=0$, whatever $\beta$;
Section~14.4 of \cite{ES17} records the same fact for non-orthogonal fields. At fixed $E_0$ the first term varies by more than a factor $10$
between the electric and the magnetic canonical field and vanishes at the magnetic end; at fixed $\epsilon$ it grows with $\beta$ (e7-27). The
critical magnetic field is $B_Q=m_e^2c^2/(e\hbar)=4.41\times10^9$~T. Radio pulsars sit at $0.02\,B_Q$; magnetars, at $10^{10}$--$10^{11}$~T, exceed
it by factors $2$--$23$ \cite{KaspiBeloborodov17} (e7-14). Yet their fields are magnetic, $\epsilon\approx0$, and create no pairs from the vacuum by
this mechanism, although their $E_0$ is the largest in nature. Hence the reading ``the rest energy of a field point materialises as pairs'' is
obstructed: the pair rate is not a function of $E_0$.

\emph{Pulsars.} In a corotating, force-free magnetosphere $\bE\cdot\Btext=0$ \cite{GoldreichJulian69}. Pair cascades need gaps where
$\bE\cdot\Btext\neq0$, which accelerate the primary particles \cite{RudermanSutherland75}. The curvature gamma-rays of the primaries convert into
pairs in the strong magnetic field \cite{Sturrock71}. That one-photon conversion takes its energy from the photon; the magnetic field supplies
momentum, not rest energy (inherited).

\emph{The analogy with the tower.} A magnetic field acts on momenta as a rotation, which is compact; an electric field acts as a boost, which is
non-compact (\cref{e7:prop:canon}). In the tower the compact elements that commute with the binding root conserve the Fock number, and every
element outside the centraliser of $K_n$, every boost included, creates quanta in pairs from the Fock vacuum (\cref{e2:lem:N}(iii),(iv),
\cref{e2:prop:bogo}). The analogy is interpretive. It shows that one algebraic dichotomy, compact against non-compact, decides both questions. It
does not identify the two. The Schwinger generator acts on the four-momentum of a charged particle in spacetime; the generators of
\cref{e2:lem:N} act on the interior Fock space, whose number $N$ is not the number of electrons. The pairs of \cref{e2:prop:bogo} appear at first
order in the rapidity, with no threshold and no exponent, as in a change of reference vacuum, whereas the Schwinger rate is non-perturbative,
with an exponent fixed by $e$ and $m_e$, neither of which the arena contains.

\subsection{Several particles in three times}
\label{e7:sec:sixmomenta}

The six-momentum $P=\sum_i(\mathsf E_i/c)\,e_i+\sum_ap_a\,e_{a+3}$ has three temporal components (Proposition~2.10 of \cite{RE}). Set $c=1$.

\begin{proposition}[Two six-momenta at mutual rest]
\label{e7:prop:sixmass}
Let $u_1,u_2\in V^+$ be unit vectors with $\eta(u_1,u_2)=\cos\vartheta$ and $P=m_1u_1+m_2u_2$. Then
\[
M^2:=\eta(P,P)=m_1^2+m_2^2+2m_1m_2\cos\vartheta ,
\]
so $M=m_1+m_2$ at $\vartheta=0$, $M=(m_1^2+m_2^2)^{1/2}$ at $\vartheta=\pi/2$ and $M=|m_1-m_2|$ at $\vartheta=\pi$. In $\Cl(3,3)$, $P^2=M^2\id$.
The graded rest Hamiltonian $H_g=m_1E_{(1,1)}+m_2E_{(2,1)}$ of the tower is the case of two orthogonal unit vectors of $\R^{6,6}$, $H_g^2=(m_1^2+m_2^2)\id$;
the ungraded $H_u=m_1\,e_1\otimes\id+m_2\,\id\otimes e_1$ has spectrum $\{\pm m_1\pm m_2\}$ and is not a vector.
\end{proposition}

\begin{proof}
Expand $\eta(P,P)$; $V^+$ is positive definite. Vectors of a Clifford algebra square to their norm. The rest is \cref{e1:prop:rest}(i),(iii).
\end{proof}

Fixed (e7-15: symbolic; $P^2$ exactly at $\cos\vartheta=\tfrac35$; the tower exactly at $n=2$). The reading is this: the graded composite of
\cref{e1:prop:rest}, with masses in quadrature (Proposition~8.11(iv) of \cite{RE}), is ``two masses with orthogonal clocks'', and the additive
ungraded pair is ``two masses with one clock''. Its precise sense is limited. In $\R^{6,6}$ the temporal triples of the two copies are orthogonal,
so $\vartheta=\pi/2$ is forced between vectors of different copies; inside one copy $\vartheta$ is free. The ungraded operator agrees with the
case $\vartheta=0$ only in its largest eigenvalue $m_1+m_2$. Its additivity is that of the eigenvalues of commuting operators, one mass shell per
particle, not that of a vector sum, which has two eigenvalues, not four. The identification is interpretive.

\begin{proposition}[Uphill decays and one clock]
\label{e7:prop:uphill}
\begin{enumerate}[label=(\roman*),leftmargin=2.2em]
\item In signature $(1,3)$, future timelike $P_1,P_2$ satisfy $\eta(P_1,P_2)\ge m_1m_2$, so $\eta(P_1+P_2,P_1+P_2)\ge(m_1+m_2)^2$: a particle
decays only into lighter ones in total.
\item In signature $(3,3)$ this fails. $P=(10,0,0;0,0,0)$ splits as $P_1=(5,5,0;1,0,0)$ plus $P_2=(5,-5,0;-1,0,0)$, with masses
$10\to7+7=14$ and $\eta(P_1,P_2)=1<49$, every clock component $\mathsf E_1$ positive.
\item The clock $e_1$ reads $m_4^2=\mathsf E_1^2-|\bp|^2=24$ for each product, and $2\sqrt{24}<10$: read by the clock the decay is downhill.
\item If all temporal parts are parallel to one unit $u\in V^+$, the momenta lie in $\R u\oplus V^-$, of signature $(1,3)$, and (i) holds.
\end{enumerate}
\end{proposition}

\begin{proof}
(i) $\eta(P_1,P_2)=E_1E_2-\bp_1\cdot\bp_2\ge E_1E_2-|\bp_1||\bp_2|\ge m_1m_2$, the last step since
$(E_1E_2-|\bp_1||\bp_2|)^2-m_1^2m_2^2=(E_1|\bp_2|-E_2|\bp_1|)^2\ge0$ with $E_1E_2\ge|\bp_1||\bp_2|$. (ii)--(iii) Integer arithmetic. (iv) The Gram
matrix of $(u,e_4,e_5,e_6)$ is $\operatorname{diag}(1,-1,-1,-1)$.
\end{proof}

(i) is inherited; (ii) is the observation that with more than one time a particle can decay into heavier ones \cite{Dorling70}, recalled in
\cite{Tegmark97}, whose main argument is the loss of well-posed prediction; an ultrahyperbolic Cauchy problem can be made well posed by a nonlocal
constraint \cite{CraigWeinstein09}, and Example~27.3 of \cite{ES17} shows that compact surplus times give tachyons. All are checked exactly or,
for (iv), with no violation among $20000$ random pairs, against $12639$ violations without a common clock (e7-16, e7-17). Item~(iii) locates the
paradox. ``Uphill'' is uphill in the six-dimensional mass $\mu$ and downhill in the mass $m_4$ read by a clock, since $m_4(n)\le\mu$ with equality iff
the temporal momentum is parallel to the clock (Proposition~3.13 of \cite{TR}; Proposition~2.10 of \cite{RE}).

\emph{Constraint.} Adopt the reading in which particles are labelled by $\mu$ and observed through one clock. Then mass additivity bounds the
temporal angles between the constituents of ordinary matter. At mutual rest $M=m_1+m_2-\mu_{12}\vartheta^2/2+O(\vartheta^4)$, with
$\mu_{12}=m_1m_2/(m_1+m_2)$; the fractional defect is $(m_1m_2/M_0^2)\vartheta^2/2$ with $M_0=m_1+m_2$, which is $\vartheta^2/8$ at equal masses,
where $M=2m\cos(\vartheta/2)$ exactly. Rainville \emph{et al.} compared mass differences with $\gamma$-ray energies and found
$1-\Delta mc^2/E=(-1.4\pm4.4)\times10^{-7}$ \cite{Rainville05}. For neutron capture on $^{28}$Si, $S_n=8.4736$~MeV from the AME2020 masses
\cite{AME2020}, so an unexplained defect is at most $3.7$~eV at $1\sigma$, a fraction $1.4\times10^{-10}$ of $M_0$. This gives
$\vartheta<9.1\times10^{-5}$~rad between the neutron and the nucleus, and $\vartheta<3.3\times10^{-5}$~rad for an equal-mass system tested at the
same fractional precision (e7-18). Proton stability enters differently. With off-clock final states accessible, the baryon-number-conserving
channel $p\to n\,e^+\nu_e$, closed in one time since $m_n>m_p$, opens kinematically by \cref{e7:prop:uphill}(ii), with a non-compact phase space. The
proton mean life exceeds $10^{29}$ years in every mode searched and $\tau/B(p\to e^+\pi^0)>2.4\times10^{34}$ years \cite{PDG,Takenaka20}. A weak
decay with MeV energy release has a lifetime of minutes unless it is forbidden, so this observation asks for a selection, one clock for all matter,
not for a small angle. The arena does not compute the rate, since it has no amplitudes. This is the classical content of the clock selection
(Proposition~3.5 of \cite{ES17}; the reduced shell of Theorem~30.1 of \cite{ES17}, Proposition~3.10 of \cite{TR}). Classes: (i) and the
kinematics of (ii) inherited; the bounds are constraints, conditional on the reading.

\subsection{\texorpdfstring{The rate $\kappa=c/\ell$ and the sky}{The rate kappa = c/l and the sky}}
\label{e7:sec:rate}

The arena fixes no $\ell$ (Theorem~4.12 of \cite{RE}). Three choices give (e7-19):
\begin{center}
\small
\begin{tabular}{@{}lccc@{}}
\toprule
$\ell$ & $\ell$ (m) & $\kappa=c/\ell$ (s$^{-1}$) & $\hbar\kappa$\\
\midrule
reduced Compton wavelength $\hbar/(m_ec)$ & $3.86\times10^{-13}$ & $7.76\times10^{20}$ & $511.0$~keV $=m_ec^2$\\
Planck length $(\hbar G/c^3)^{1/2}$ & $1.62\times10^{-35}$ & $1.85\times10^{43}$ & $1.22\times10^{19}$~GeV\\
Hubble radius $c/H_0$ \cite{Planck18VI} & $1.37\times10^{26}$ & $2.18\times10^{-18}$ & $1.44\times10^{-33}$~eV\\
\bottomrule
\end{tabular}
\end{center}
The first is the identification $mc^2=\hbar\kappa$ of \cite{RE}, for which $\ell=\hbar/(mc)$ is the only length (Theorem~4.12(ii) of \cite{RE}). The
last puts $\hbar\kappa$ at the scale of ultralight fields. Each is a reading with no support from the arena.

\begin{proposition}[Duality on a null field is a rotation about its axis]
\label{e7:prop:dualrot}
Let $z=\bB+i\bE\neq0$ be null and $\mathbf n=\bB\times\bE/|\bB\times\bE|$, which equals $\bE\times\Btext/|\bE\times\Btext|$, the direction of
propagation. Then $\mathbf n\times z=-iz$, $\mathbf n\cdot z=0$, and the duality rotation $e^{\vartheta G}$, $z\mapsto e^{i\vartheta}z$, acts as the
rotation $R_{\mathbf n}(-\vartheta)$ on $\bB$ and on $\bE$ separately. The plane wave $z=w\,f(\mathbf n\cdot\mathbf x-t)$ with $\mathbf n\times w=-iw$
solves $\partial_tz=-i\curl z$ (Proposition~12.10 of \cite{ES17}). For $q\neq0$ and $\vartheta\notin\pi\Z$ no rotation acts as $e^{\vartheta G}$.
\end{proposition}

\begin{proof}
$\mathbf n\times z=-iz$ is Theorem~3.4(ii) of \cite{RE}, from $\bB\cdot\bE=0$ and $|\bB|=|\bE|$; dotting with $\mathbf n$ gives $\mathbf n\cdot z=0$. Then
$e^{i\vartheta}z=\cos\vartheta\,z+\sin\vartheta\,iz=\cos\vartheta\,z-\sin\vartheta\,\mathbf n\times z=R_{\mathbf n}(-\vartheta)z$ by Rodrigues' formula,
and $R_{\mathbf n}(-\vartheta)$ is real. For the wave, $\curl z=\mathbf n\times w\,f'=-iw\,f'$ and $\partial_tz=-w\,f'$. Rotations fix $q$, while
$e^{\vartheta G}$ sends $q$ to $e^{2i\vartheta}q$.
\end{proof}

Fixed; checked on $50$ null fields at three angles to $2\times10^{-15}$, the wave equation exactly (e7-20). It is the identity of helicity and
duality on null fields \cite{Calkin65,PenroseRindler84}, Example~13.11 of \cite{ES17}, pulled back to the frame. Physically, a duality angle
$\vartheta$ turns the plane of linear polarisation of a freely propagating wave about its line of sight by $\vartheta$, the same for all
frequencies, and only multiplies a circularly polarised wave by a phase.

\emph{Constraint, and an obstructed reading.} Suppose the interior flow $X(\tau)=e^{\mp\kappa G\tau}X_0$ (Proposition~4.9(ii) of \cite{RE}) acted on
freely propagating radiation, with $\tau$ the cosmic time. By \cref{e7:prop:dualrot} linearly polarised light would arrive with its plane turned by
$\kappa t$. The CMB polarisation gives a uniform rotation $\beta=0.35^\circ\pm0.14^\circ$ (68\%) \cite{MinamiKomatsu20}; distant radio galaxies bound
a uniform rotation over cosmological distances to a few degrees \cite{CarrollFieldJackiw90}. With $|\beta|<0.63^\circ$ ($|$central value$|+2\sigma$)
over $t_0=13.79$~Gyr \cite{Planck18VI}, the light having travelled for nearly all of it, $\kappa<2.5\times10^{-20}\,{\rm s^{-1}}=1.2\times10^{-2}H_0$,
that is $\hbar\kappa<1.7\times10^{-35}$~eV. This is $10^{40.5}$ below $m_ec^2/\hbar=7.8\times10^{20}\,{\rm s^{-1}}$ (e7-21). Under $mc^2=\hbar\kappa$,
for the electron or any particle heavier than $1.7\times10^{-35}$~eV, the interior flow cannot act on propagating radiation: that reading is
obstructed. Even the Hubble rate is excluded by a factor $86$.

\emph{The converse reading.} A cosmological duality angle of light is cosmic birefringence, as produced by a Chern--Simons or axion-like term
\cite{CarrollFieldJackiw90} (Section~14.4 of \cite{ES17}). In the closure its potential vanishes on null fields. $U(\theta)=(R_0/2m)\cos\psi$ with
$R_0=2|\nu|E_0(\pi_0)=0$, so the duality angle of radiation is an exact flat direction (Proposition~4.3 and Corollary~4.5(i) of \cite{RE}; restated,
e7-22). The closure therefore neither forbids nor predicts a uniform rotation of polarisation. Its rate is a free datum, open, bounded as above
if constant in time. A measured $\beta\neq0$, now at $2.4\sigma$, would fix that datum and not test the arena.

\subsection{Time-odd invariants, clocks and the constancy of masses}
\label{e7:sec:timeodd}

\emph{Three manifolds and three generations.} For three manifolds with rates $\kappa_k$ and couplings $B_{kl}$ the plaquette is
$H_{12}H_{23}H_{31}=iJ$ with $J=B_{12}B_{23}B_{31}$, $\det[H,K]=2iJ(\kappa_1-\kappa_2)(\kappa_1-\kappa_3)(\kappa_2-\kappa_3)$, and interior time
reversal sends $J$ to $-J$, which no copy reflection undoes at $n=3$ (Theorem~9.5(ii),(iv) of \cite{RE}; restated, e7-23). The Jarlskog invariant of
quark mixing has the same shape: a rephasing-invariant plaquette of the mixing matrix, nonzero only for three generations, with
$\det[M_uM_u^\dagger,M_dM_d^\dagger]$ equal to $2iJ$ times the two Vandermonde products of squared masses \cite{Jarlskog85,KobayashiMaskawa73}. The
parallel is a reading. It says that in both cases a time-odd invariant of a Hermitian $3\times3$ structure needs three labels, and that the
commutator determinant factorises through it. It does not establish CP violation in particle physics: the arena has no baryons, no C, no
CP, and its $J$ is a property of the interior $b$-field between manifolds, with a free magnitude. Sakharov's conditions for a baryon asymmetry,
baryon-number violation, C and CP violation and departure from equilibrium \cite{Sakharov67}, are not touched. At most $J$ is a candidate source
of time-odd effects, and only under a reading that identifies interior time reversal with physical time reversal.

\emph{A temporal $b$-field on the clocks of matter.} A temporal rotation turns the temporal momentum at fixed six-dimensional mass, and the mass
read by a fixed clock oscillates (\cref{e6:prop:clock}). At rest, for a libration $\phi=\theta_0\sin\Omega t$ in the plane $(e_1,e_2)$ read by
$e_1$, $m_4=\mu\cos\phi$ and $m_4/\mu=1-\theta_0^2/4+(\theta_0^2/4)\cos2\Omega t+O(\theta_0^4)$. For a uniform rotation at rate $\nu$ about an axis at
angle $\alpha$ from the clock, with the temporal energy initially on the clock, $m_4/\mu=1-\sin^2\alpha\,(1-\cos\nu t)$ (fixed, e7-24). Both are
$m_4/\mu=1-A+A\cos\omega t$. A rotation common to all matter cancels in dimensionless mass ratios, so the bounds concern differential
oscillations; all are conditional on the reading (e7-25).
\begin{itemize}
\item \emph{Quasar absorption.} Methanol lines at $z=0.89$ give $|\Delta\mu_{pe}/\mu_{pe}|\lesssim4\times10^{-7}$ ($2\sigma$) for
$\mu_{pe}=m_p/m_e$, over a look-back time of about $7.5$~Gyr \cite{Kanekar15}. For generic phases this gives $A<4\times10^{-7}$, so
$\theta_0<1.3\times10^{-3}$~rad for the libration and $\alpha<6.3\times10^{-4}$~rad for the uniform rotation. One pair of epochs cannot exclude a
coincidence of phases.
\item \emph{Atomic clocks.} Yb$^+$ against Cs gives $\dot\mu_{pe}/\mu_{pe}=(-0.5\pm1.6)\times10^{-16}\,{\rm yr^{-1}}$ \cite{Huntemann14}, consistent with
$(0.2\pm1.1)\times10^{-16}\,{\rm yr^{-1}}$ \cite{Godun14}. At $2\sigma$, $A\omega<3.7\times10^{-16}\,{\rm yr^{-1}}$ for periods long compared with the
years of comparison: $A<2.6\times10^{-6}$ for $\omega$ equal to the inverse age, and $A<1.8\times10^{-16}$ for $\omega=1\,{\rm yr^{-1}}$.
\item \emph{Equivalence principle.} Add the hypothesis that inertia is $m_4$ read by the clock while gravity couples to the undivided $\mu$. Then
$m_g/m_i=1/\cos\vartheta\approx1+\vartheta^2/2$, and the E\"otv\"os parameter of two bodies is $\eta\approx\half(\vartheta_A^2-\vartheta_B^2)$.
MICROSCOPE measured $\eta({\rm Ti},{\rm Pt})=[-1.5\pm2.3({\rm stat})\pm1.5({\rm syst})]\times10^{-15}$ \cite{Touboul22}, so $|\eta|<7.0\times10^{-15}$
at $2\sigma$, $|\vartheta_A^2-\vartheta_B^2|<1.4\times10^{-14}$, and $\vartheta<1.2\times10^{-7}$~rad if one body is aligned. If inertia and gravity
both couple to $m_4$, the test says nothing.
\end{itemize}

\subsection{Summary and answer}
\label{e7:sec:summary}

\begin{table}[ht]
\centering
\footnotesize
\begin{tabular}{@{}p{3.3cm}p{3.3cm}p{2.5cm}p{1.6cm}p{3.6cm}@{}}
\toprule
finding & classical counterpart & class & check & what an observation could decide\\
\midrule
$u_0$ of a field point (RE~2.8) & energy density where $\bE\parallel\Btext$; minimum over frames & restated; inherited & e7-2 & nothing (identity)\\
traceless stress, tension and pressure $u_0$ (\cref{e7:prop:stress}) & radiation, $w=\tfrac13$; von Laue & inherited & e7-3, e7-4 & nothing\\
$\rho_0$ as dark matter (\cref{e7:cor:dm}) & pressureless dark matter & obstructed & e7-26 & already decided: dark matter has $w\approx0$\\
$u_0$ of astronomical fields (\cref{e7:tab:fields}) & magnetic energy densities; ISM equipartition & inherited arithmetic & e7-5, e7-6 & nothing for the arena\\
$E_0/E_6$ has density $2t$ (\cref{e7:prop:gauss}) & Gaussian thermal field & fixed; model inherited & e7-7 to e7-9 & joint statistics of $\bE,\Btext$ at a point test the model\\
CMB ``locally massive'', $\langle u_0\rangle=\tfrac23u$ & photon gas, $Mc^2=U$ & reading & e7-10 & nothing gravitational\\
$M_{\rm pair}$ of two null fields (RE~8.6) & Breit--Wheeler, gamma-ray horizon & restated; same zero set; thresholds inherited & e7-11, e7-12 & horizon tests QED and background light\\
canonical fields, $E_0=\half(\epsilon^2+\beta^2)$ (\cref{e7:prop:canon}) & Schwinger rate, $E_S$, $B_Q$ & inherited & e7-13, e7-14 & strong-field QED\\
rest energy materialises as pairs & magnetars above $B_Q$ & obstructed & e7-27 & already decided (Schwinger)\\
compact/boost and pairs (\cref{e2:lem:N}) & magnetic against electric fields & interpretive (analogy) & e7-13 & nothing\\
$M^2=m_1^2+m_2^2+2m_1m_2\cos\vartheta$ (\cref{e7:prop:sixmass}) & graded quadrature, additive copies & fixed; identification interpretive & e7-15 & nothing\\
uphill decays (\cref{e7:prop:uphill}) & Dorling's instability & inherited; example fixed & e7-16, e7-17 & proton stability: one clock\\
one clock for matter & mass additivity, $E=mc^2$ test & constraint (conditional) & e7-18 & $\vartheta<9\times10^{-5}$~rad; better tests tighten\\
$\kappa=c/\ell$, three lengths & Compton, Planck, Hubble & interpretive & e7-19 & nothing until $\ell$ is fixed\\
duality $=$ rotation about $\mathbf n$ (\cref{e7:prop:dualrot}) & helicity, polarisation rotation & fixed; inherited & e7-20 & nothing\\
interior flow on light at $\kappa=m_ec^2/\hbar$ & cosmic birefringence & obstructed; constraint $\kappa<2.5\times10^{-20}$~s$^{-1}$ & e7-21 & CMB polarisation\\
flat duality angle of radiation (RE~4.5) & axion-like birefringence & open; restated & e7-22 & $\beta$ fixes a datum, tests nothing\\
three-manifold $J$ (RE~9.5) & Jarlskog, Sakharov & reading & e7-23 & nothing, unless manifolds become families\\
oscillating $m_4$ (\cref{e6:prop:clock}) & constancy of $m_p/m_e$, equivalence principle & fixed; bounds conditional & e7-24, e7-25 & quasars, clocks, EP tests\\
\bottomrule
\end{tabular}
\caption{Claims of \cref{e7:sec:main}: the finding, its classical counterpart, its class, its check (\texttt{PASS} lines of
\texttt{verify/e7\_classical.py}, e7-1 being the realisation) and what an observation could decide. ``RE~$x$'' is Proposition, Theorem or
Corollary~$x$ of \cite{RE}.}
\label{e7:tab:claims}
\end{table}

\emph{Answer.} Classical or astronomical content belongs to the inherited parts: the rest energy of a field point and its stress, the
pointwise rest fraction $\tfrac23$ of thermal radiation, pair kinematics and the gamma-ray horizon, the canonical fields and the Schwinger rate,
the kinematics of several times, and the identity of duality and helicity on light. There the arena agrees with ordinary physics and adds
nothing measurable. Excluded as readings are field rest energy as dark matter, the materialisation of rest energy into pairs, the interior flow
acting on light at the rate $m_ec^2/\hbar$, and six-dimensional kinematics without a common clock. Unconstrained, or open, are $\ell$ and with it
$\kappa$, the origin of a cosmic birefringence angle, the magnitude and meaning of the three-manifold invariant, and the readings that identify
the interior with particles. Two structures are missing before the arena can make an astronomical or cosmological prediction. One is a fixed
length $\ell$, equivalently a mass through $mc^2=\hbar\kappa$. The other is a coupling of the interior to spacetime: on which fields the interior
flow acts, how the temporal part of the six-momentum couples to clocks and to gravity, and which clock ordinary matter shares. With them, the
constraints above become tests: the angles of \cref{e7:sec:sixmomenta,e7:sec:timeodd} and the rate of \cref{e7:sec:rate} would be predicted
rather than bounded.
